\documentclass[UTF8]{ctexart}
\usepackage{xeCJK}
\usepackage{geometry}
\geometry{left=2cm,right=2cm,top=2.5cm,bottom=2.5cm}
\setlength{\headheight}{12.7pt}

\usepackage{titlesec}
\titleformat{\section}
  {\normalfont\large\bfseries}
  {\thesection}
  {1em}
  {}
\titlespacing*{\section}
  {0pt}
  {3.5ex plus 1ex minus .2ex}
  {2.3ex plus .2ex}

\usepackage{amsmath}
\usepackage{amssymb}
\usepackage{amsthm}
\usepackage{enumitem}
\setlist[enumerate]{itemsep=0pt,topsep=0pt,parsep=0pt,leftmargin=0pt,itemindent=2.5em}
\setlist[itemize]{itemsep=0pt,topsep=0pt,parsep=0pt,leftmargin=0pt,itemindent=2.5em}
\usepackage{fancyhdr}
\pagestyle{fancy}
\fancyhf{}
\usepackage{graphicx}
\usepackage{hyperref}
\usepackage{indentfirst}
\usepackage{lmodern}


\begin{document}
\setlength{\abovedisplayskip}{1pt}
\setlength{\belowdisplayskip}{1pt}

\section{等价关系}

等价关系是一种特殊的\href{02 二元关系#nameddest=sec:定义1.1}{二元关系}.
这里把$(x,y)\in\,\sim$记为$x\sim y$.

\vspace{10pt}

\hypertarget{sec:定义1.1}{}
\noindent\textbf{定义1.1 (等价关系)}

给定集合$a$和$a$上的二元关系$\sim$, 如果以下三条成立:
\begin{enumerate}
    \item 自反性: $\forall x\in a:\,x\sim x$,
    \item 对称性: $\forall x,y\in a:\,x\sim y\to y\sim x$,
    \item 传递性: $\forall x,y,z\in a:\,x\sim y\land y\sim z\to x\sim z$,
\end{enumerate}
那么称$\sim$是$a$上的一个\textbf{等价关系}.
也称$(a,\sim)$是一个\textbf{等价关系结构}.

\vspace{10pt}

\hypertarget{sec:定义1.2}{}
\noindent\textbf{定义1.2 (等价类、商集)}

给定等价关系结构$(a,\sim)$. 对任意的$x\in a$, 定义
\[
    [x]\overset{\mathrm{def}}{=}\{y\in a\mid x\sim y\},
\]
称之为$x$在$\sim$中的\textbf{等价类}. 进一步, 定义
\[
    ^{\displaystyle a}\!/_{\displaystyle\sim}\overset{\mathrm{def}}{=}
    \{\,[x]\mid x\in a\},
\]
称之为$\sim$的\textbf{商集}.

\vspace{10pt}

一个等价关系等同于集合的一个划分.

\vspace{10pt}

\hypertarget{sec:定义1.3}{}
\noindent\textbf{定义1.3 (集合的划分)}

给定集合$a$和它的子集族$S\subset\mathcal{P}(a)$. 如果:
\begin{enumerate}
    \item $\varnothing\notin S$,
    \item $\forall s,s'\in S:\,s\neq s'\to s\cap s'=\varnothing$,
    \item $\displaystyle\bigcup S=a$, 即$\forall x\in a\,\exists s\in S:x\in s$,
\end{enumerate}
那么称$S$是$a$的一个\textbf{划分}.

\vspace{10pt}

\hypertarget{sec:定理1.1}{}
\noindent\textbf{定理1.1}

给定集合$a$, 记其上所有等价关系的集合为$\mathcal{E}$,
记其所有划分的集合为$\mathcal{F}$, 则有典范\href{03 映射#nameddest=sec:定义2.2}{双射}
\[
    \varphi:\mathcal{E}\to\mathcal{F},\quad\sim\,\,\mapsto
    \,\!^{\displaystyle a}\!/_{\displaystyle\sim}\,;\quad\quad\quad
    \varphi^{-1}:\mathcal{F}\to\mathcal{E},\quad S\mapsto
    \{(x,y)\in a^2\mid\exists s\in S:x,y\in s\}.
\]

\noindent{\kaishu 证明.}

一方面, 任给等价关系$\sim$, 商集$\,\!^{\displaystyle a}\!/_{\displaystyle\sim}$
是$a$的划分, 因为:
\begin{enumerate}
    \item 等价类都不是空集,
    \item 两个不同的等价类$[x],[y]$的交集是空集, 否则有$z\in[x]\cap[y]$,
    则$x\sim z\sim y$, 显然$[x]=[y]$, 矛盾,
    \item 对任意的$x\in a$, 等价类$[x]$含$x$,
\end{enumerate}
并且对任意的$x,y\in a$, $x\sim y$当且仅当它们同属一个等价类.

另一方面, 任给划分$S$, $x\sim y\overset{\mathrm{def}}{\leftrightarrow}
\exists s\in S:x,y\in s$是$a$上的等价关系, 因为:
\begin{enumerate}
    \item 对任意的$x\in a$, 存在$s$使得$x\in s$, 即$x\sim x$,
    \item 对任意的$x,y\in a$, 如果$x,y$同属$s$, 显然$y,x$同属$s$,
    \item 对任意的$x,y\in a$, 如果$x,y$同属$s$, 并且$y,z$同属$s'$,
    那么$s=s'$并且$x,z$同属于它,
\end{enumerate}
并且每个$s\in S$恰好是一个等价类, 即$S$恰好是商集. \hfill $\qedsymbol$





\newpage

\hypertarget{sec:定义1.4}{}
\noindent\textbf{定义1.4 (选代表函数)}

给定等价关系结构$(a,\sim)$和映射
$C:\,\!^{\displaystyle a}\!/_{\displaystyle\sim}\to a$. 如果
\[
    \forall s\in\,\!^{\displaystyle a}\!/_{\displaystyle\sim}:\quad C(s)\in s,
\]
就称$C$是$(a,\sim)$的一个\textbf{选代表函数}, 称$C(s)$为等价类$s$的代表.

\vspace{10pt}

\hypertarget{sec:定理1.2}{}
\noindent\textbf{定理1.2}

任意的等价关系结构都有选代表函数, 选代表函数一定是单射.

\noindent{\kaishu 证明.}

商集上的选择函数就是选代表函数.
如果$C(s)=C(s')$, 那么它同时属于$s$和$s'$, 从而$s=s'$. \hfill $\qedsymbol$

\vspace{10pt}

最后, 数学中经常对等价类定义映射或命题, 下面描述这种商操作的合理性.

\vspace{10pt}

\hypertarget{sec:定理1.3}{}
\noindent\textbf{定理1.3 (商映射)}

给定映射$f:a\to b$和$a$上的等价关系$\sim$. 如果
\[
    \forall x,y\in a:\quad x\sim y\to f(x)=f(y),
\]
那么存在唯一的映射$F:\,\!^{\displaystyle a}\!/_{\displaystyle\sim}\to b$使得
\[
    \forall s\in\,\!^{\displaystyle a}\!/_{\displaystyle\sim}\,\,
    \forall x\in s:F(s)=f(x),\\[1pt]
\]
称之为映射$f$在$\sim$下的\textbf{商映射}.

\noindent{\kaishu 证明.}

\noindent{\kaishu 存在性.}

取一个选代表函数$C$, 令$F=f\circ C$, 则对任意的等价类$s$和$x\in s$, 有
\[
    C(s)\in s\implies C(s)\sim x\implies F(s)=f(C(s))=f(x).
\]

\noindent{\kaishu 唯一性.} 显然. \hfill $\qedsymbol$

\vspace{10pt}

\hypertarget{sec:定理1.4}{}
\noindent\textbf{定理1.4 (商命题)}

给定命题$p(x)$和$a$上的等价关系$\sim$. 如果
\[
    \forall x,y\in a:\quad x\sim y\to(p(x)\leftrightarrow p(y)),
\]
那么存在唯一的命题$P(s)$使得
\[
    \forall s\in\,\!^{\displaystyle a}\!/_{\displaystyle\sim}\,\,
    \forall x\in s:(P(s)\leftrightarrow p(x)),\\[1pt]
\]
称之为命题$p$在$\sim$下的\textbf{商命题}.

\noindent{\kaishu 证明.}

\noindent{\kaishu 存在性.}

取一个选代表函数$C$, 定义$P(s)\overset{\mathrm{def}}{\leftrightarrow}p(C(s))$,
则对任意的等价类$s$和$x\in s$, 有
\[
    C(s)\in s\implies C(s)\sim x\implies p(C(s))\leftrightarrow p(x)
    \implies P(s)\leftrightarrow p(x).
\]

\noindent{\kaishu 唯一性.} 显然. \hfill $\qedsymbol$

\end{document}